Give the small dictionary of statistics to algebraic geometry

Probability/Statistics
Probability distribution
Statistical model
Discrete exponential family
Conditional inference
Maximum likelihood estimation
Model selection
Multivariate Gaussian model
Phylogenetic model
MAP estimates

Algebra/Geometry
Point
(Semi)Algebraic set
Toric variety
Lattice points in polytopes
Polynomial optimization
Geometry of singularities
Spectrahedral geometry
Tensor networks
Tropical geometry


% Probabiliity primer

State how all multivariate normal distributions can be obtained from affine transforms of the standard mv. normal


We say that a random vector \( X = (X_1, \ldots, X_m) \) is distributed according to the standard multivariate normal distribution 
\( \mathcal{N}_m(0, I) \) if its components \( X_1, \ldots, X_k \) are independent and identically distributed as \( X_i \sim \mathcal{N}(0,1) \). 
Thus, the distribution \( \mathcal{N}_m(0, I) \) has the density function
\[
\varphi_{0,I}(x) = \prod_{i=1}^m \frac{1}{\sqrt{2\pi}} \exp\left\{ -\frac{1}{2} x_i^2 \right\}
= \frac{1}{(2\pi)^{m/2}} \exp\left\{ -\frac{1}{2} x^T x \right\}.
\]

Let \( \Lambda \in \mathbb{R}^{m \times m} \) be a matrix of full rank and let \( \mu \in \mathbb{R}^m \). Suppose \( X \sim \mathcal{N}_m(0,I) \) 
and define \( Y = \Lambda X + \mu \). Integratin by substitution implies, that the random vector \( Y \) has the density function
\[
\varphi_Y(y) = \frac{1}{(2\pi)^{m/2}} \exp\left\{ -\frac{1}{2} (y-\mu)^T \Lambda^{-T} \Lambda^{-1} (y-\mu) \right\} |\Lambda^{-1}|,
\]
where \(|\cdot|\) denotes the determinant, and \(\Lambda^{-T}\) denotes the transpose of the inverse. 
Define \(\Sigma = \Lambda \Lambda^T\). Then the density \(\varphi_Y\) is seen to depend only on \(\mu\) and \(\Sigma\) and can be rewritten as
\[
\varphi_{\mu,\Sigma}(y) = \frac{1}{(2\pi)^{m/2} |\Sigma|^{1/2}} \exp\left\{ -\frac{1}{2} (y-\mu)^T \Sigma^{-1} (y-\mu) \right\}.
\]
Note that since the standard multivariate normal random vector has expectation \( 0 \) and covariance matrix \( I \), it follows 
that \( Y \) has expectation \( \Lambda 0 + \mu = \mu \) and covariance matrix \( \Lambda \Lambda^T = \Sigma \).


State the conditionals of multivariate normal distributions and distribution of subindices

Theorem. Let \( X \sim \mathcal{N}_m(\mu, \Sigma) \) and let \( A, B \subseteq [m] \) be two disjoint index sets.

(i) The marginal distribution of \( X \) over \( B \), that is, the distribution of the subvector \( X_B = (X_i \mid i \in B) \), 
is multivariate normal; namely,
\[
X_B \sim \mathcal{N}_{\#B}(\mu_B, \Sigma_{B,B}),
\]
where \( \mu_B \) is the respective subvector of the mean vector \( \mu \) and \( \Sigma_{B,B} \) is the corresponding principal submatrix of 
\( \Sigma \).

(ii) For every \( x_B \in \mathbb{R}^B \), the conditional distribution of \( X_A \) given \( X_B = x_B \) is the multivariate normal distribution
\[
\mathcal{N}_{\#A}\left( \mu_A + \Sigma_{A,B}\Sigma_{B,B}^{-1}(x_B - \mu_B),\, \Sigma_{A,A} - \Sigma_{A,B}\Sigma_{B,B}^{-1}\Sigma_{B,A} \right).
\]

The matrix \( \Sigma_{A,B}\Sigma_{B,B}^{-1} \) appearing in the theorem is the matrix of regression coefficients. The positive definite matrix
\[
\Sigma_{AA.B} = \Sigma_{A,A} - \Sigma_{A,B}\Sigma_{B,B}^{-1}\Sigma_{B,A}
\]
is the conditional covariance matrix. We also note that \( (X_A \mid X_B = x_B) \) is not well-defined as a random vector. It is merely a way of denoting the conditional distribution of \( X_A \) given \( X_B = x_B \).


% Algebra preliminaries

State the elimination theorem and the vanishing ideal of the graph of a rational map


Recall the Elimination theorem:
Let \(\pi : \mathbb{K}^{r_1 + r_2} \to \mathbb{K}^{r_1}\) be the coordinate projection
\[
(a_1, \ldots, a_{r_1}, b_1, \ldots, b_{r_2}) \mapsto (a_1, \ldots, a_{r_1}).
\]

Let \(V \subseteq \mathbb{K}^{r_1 + r_2}\) be a variety and let
\[
I = I(V) \subseteq \mathbb{K}[p, q] := \mathbb{K}[p_1, \ldots, p_{r_1}, q_1, \ldots, q_{r_2}]
\]
be its vanishing ideal. Then the vanishing ideal of the image \( \pi(V) \) is the
ideal
\[ I(\pi(V)) = I \cap \mathbb{K}[p] \]
where RHS is also called the elimination ideal.
For practical purposes also keep in mind that for an ordering \( p_1 \prec \dots \prec q_{r_2} \),
that if \( \mathcal{G} \) is a Groebner basis of \( I \), \( \mathcal{G} \cap \mathbb{K}[p] \) is a Groebner basis
for \( I \cap \mathbb{K}[p] \).

The generalization to arbitrary rational maps (instead of projections) is:
Suppose that \(\phi : \mathbb{K}^{r_2} \to \mathbb{K}^{r_1}\) is a polynomial map (each component function is a polynomial). The \emph{graph} of the map is the set of points
\[
\Gamma_\phi = \{ (\phi(\theta), \theta) : \theta \in \mathbb{K}^{r_2} \} \subset \mathbb{K}^{r_1 + r_2}.
\]

If \(\mathbb{K}\) is an infinite field the vanishing ideal of the graph is
\[
I(\Gamma_\phi) = \langle p_1 - \phi_1(q), \ldots, p_{r_1} - \phi_{r_1}(q) \rangle.
\]

Let \(I \subset \mathbb{K}[q]\) be the vanishing ideal of a variety \(V \subset \mathbb{K}^{r_2}\). 
Let \(\phi\) be a polynomial map \(\phi : \mathbb{K}^{r_2} \to \mathbb{K}^{r_1}\). 
Then the vanishing ideal of the Zariski closure is the elimination ideal
\[
I(\phi(V))\quad =\quad (I + I(\Gamma_\phi)) \cap \mathbb{K}[p].
\]
If \( \phi : \mathbb{K}^{r_2} \to \mathbb{K}^{r_1} \) is rational with coord. functions \( \phi_i = \frac{f_i}{g_i} \) 
then
\[
    I(\phi(V)) = \left(I + \langle g_1 p_1 - f_1, \dots, g_{r_1}p_{r_1} - f_{r_1}, zg_1 \cdots g_{r_1} - 1 \rangle \right) \cap \mathbb{K}[p]
\]
The \( zg_1 \cdots g_{r_1} - 1 \) is necessary to consider only points where the denominator doesn't vanish and \( z \) is necessary
to clear denominators.


Give an application of the elimination theorem for the binomial distributions (Implicitization)

Consider the parametrization of the binomial random variable model with two trials from. 
To compute the implicit description of this model, we can start with the vanishing ideal \(I\) of the probability simplex \(\Delta_1\), 
which is the ideal
\[
I = \langle s + t - 1 \rangle \subseteq \mathbb{R}[s, t].
\]

The parametrization has the form
\[
\phi(s, t) = (s^2, 2st, t^2).
\]

We form the ideal of the graph
\[
I + I(\Gamma_\phi) = \langle s + t - 1,\, p_0 - s^2,\, p_1 - 2st,\, p_2 - t^2 \rangle \subseteq \mathbb{R}[s, t, p_0, p_1, p_2].
\]

Computing a lexicographic Gröbner basis for this ideal with \(s \succ t \succ p_0 \succ p_1 \succ p_2\) produces the vanishing ideal of the model
\[
\langle p_0 + p_1 + p_2 - 1,\, 4p_0p_2 - p_1^2 \rangle.
\]


What is the primary decomposition? State related facts.

Proposition.
1. Every irreducible ideal is primary.
2. If \( Q \) is primary, then \( \sqrt{Q} = P \) is a prime ideal, called the associated prime of \( Q \).
3. If \( Q_1 \) and \( Q_2 \) are primary, and \( \sqrt{Q_1} = \sqrt{Q_2} = P \), then \( Q_1 \cap Q_2 \) is also primary, 
   with associated prime \( P \).

Definition 
Let \( I \subseteq \mathbb{K}[p] \) be an ideal. A primary decomposition of \( I \)
is a representation \( I = \bigcap_{i=1}^r Q_i \); an intersection of finitely primary ideals. 
The decomposition is irredundant if no \( \bigcap_{j\neq i} Q_j \subseteq Q_i \) and
all the associated primes \( \sqrt{Q}_i \) are distinct.

Theorem
Let \(I\) be an ideal in a polynomial ring.
- Then \(I\) has an irredundant primary decomposition \(I = Q_1 \cap \cdots \cap Q_k\) such that the prime ideals \(\sqrt{Q_i}\) are all distinct.
- The set of prime ideals \(\mathrm{Ass}(I) := \{\sqrt{Q_i} : i \in [k]\}\) is uniquely determined by \(I\), and does not depend on the particular irredundant primary decomposition. They are called the associated primes of \(I\).
- The minimal elements of \(\mathrm{Ass}(I)\) with respect to inclusion are called minimal primes of \(I\). The primary components \(Q_i\) associated to minimal primes of \(I\) are uniquely determined by \(I\).
- The minimal primes of \(I\) are precisely the minimal elements among all prime ideals that contain the ideal \(I\).
- Associated primes of \(I\) that are not minimal are called embedded primes. The primary components \(Q_i\) associated to embedded primes of \(I\) are not uniquely determined by \(I\).
- \(\sqrt{I} = \bigcap \sqrt{Q_i}\). If the intersection is taken over the minimal primes of \(I\), this gives a primary (in fact, prime) decomposition of \(\sqrt{I}\).
- The decomposition \(V(I) = \bigcup V(\sqrt{Q_i})\), where the union is taken over the minimal primes of \(I\), is the irreducible decomposition of the variety \(V(I)\).

Note that this can be interpreted geometrically that any variety has a decomposition into irreducible subvarieties \( V = \bigcup_{i=1}^{r} V_i \),
where \( V_i \neq W_1 \cup W_2 \) (irreducible) for any proper subvarieties \( W_1, W_2 \subset V \). Also recall that a variety is irreducible
if and only if \( I(V) \) is a prime ideal.

% Conditional independence

State the conditional independence inference rules

Let \(A, B, C, D \subseteq [m]\) be pairwise disjoint subsets. Then
(i) \emph{(symmetry)} \quad \(X_A \perp\!\!\!\perp X_B \mid X_C \implies X_B \perp\!\!\!\perp X_A \mid X_C;\)
(ii) \emph{(decomposition)} \quad \(X_A \perp\!\!\!\perp X_{B \cup D} \mid X_C \implies X_A \perp\!\!\!\perp X_B \mid X_C;\)
(iii) \emph{(weak union)} \quad \(X_A \perp\!\!\!\perp X_{B \cup D} \mid X_C \implies X_A \perp\!\!\!\perp X_B \mid X_{C \cup D};\)
(iv) \emph{(contraction)} \quad \(X_A \perp\!\!\!\perp X_B \mid X_{C \cup D}\) and \(X_A \perp\!\!\!\perp X_D \mid X_C\) \(\implies X_A \perp\!\!\!\perp X_{B \cup D} \mid X_C.\)

Suppose that the join distribution satisfies \( f(x) > 0 \) for all \( x \in \Omega \) then
\[ X_A \perp\!\!\!\perp X_B \mid X_{C \cup D} \text{ and } X_A \perp\!\!\!\perp X_C \mid X_{B \cup D} \implies X_A \perp\!\!\!\perp X_{B \cup C} \mid X_D; \]

The conditional independence ideal:
Let \( X = (X_1, \ldots, X_m) \) be a vector of discrete random variables. 
Returning to the notation used in previous chapters, we let \([r_j]\) be the set of values taken by \(X_j\). 
Then \(X\) takes its values in \(\mathcal{R} = \prod_{j=1}^m [r_j]\). 
For \(A \subseteq [m]\), \(\mathcal{R}_A = \prod_{a \in A} [r_a]\). 
In this discrete setting, a conditional independence constraint translates into a system of quadratic polynomial equations in the 
joint probability distribution.

Proposition: 
If \(X\) is a discrete random vector, then the conditional independence statement \(X_A \perp\!\!\!\perp X_B \mid X_C\) holds if and only if
\[
p_{i_A, i_B, i_C, +} \cdot p_{j_A, j_B, i_C, +}
-
p_{i_A, j_B, i_C, +}
\cdot
p_{j_A, i_B, i_C, +}
=
0
\]
for all \(i_A, j_A \in \mathcal{R}_A\), \(i_B, j_B \in \mathcal{R}_B\), and \(i_C \in \mathcal{R}_C\).
The ideal \( I_{X_A \perp\!\!\!\perp X_B \mid X_C} \)independence generated by these quadr. polynomials is called the conditional independence ideal.


State the conditional independence condition for multivariate normals and relate it to primary decompositions

Proposition
The conditional independence statement \( X_A \perp\!\!\!\perp X_B \mid X_C \) holds for a multivariate normal random vector 
\( X \sim \mathcal{N}(\mu, \Sigma) \) if and only if the submatrix \( \Sigma_{A \cup C,\, B \cup C} \) of the covariance matrix \( \Sigma \) 
has rank \(\# C\).

Definition:
Fix pairwise disjoint subsets \(A\), \(B\), and \(C\) of \([m]\). The Gaussian conditional independence ideal 
\(J_{A \perp\!\!\!\perp B \mid C}\) is the following ideal in \(\mathbb{R}[ \sigma_{ij},\, 1 \leq i \leq j \leq m ]\):
\[
J_{A \perp\!\!\!\perp B \mid C} = \langle (\#C + 1) \times (\#C + 1)\ \text{minors of}\ \Sigma_{A \cup C,\, B \cup C} \rangle.
\]


Similarly as with the discrete case, 
we can take a collection \(\mathcal{C}\) of conditional independence constraints, and get a conditional independence ideal
\[
J_{\mathcal{C}} = \sum_{A \perp\!\!\!\perp B \mid C \in \mathcal{C}} J_{A \perp\!\!\!\perp B \mid C}.
\]
The resulting Gaussian conditional independence model is a subset of the cone \(PD_m\) of \(m \times m\) symmetric positive definite matrices:
\[
\mathcal{M}_{\mathcal{C}} := V(J_{\mathcal{C}}) \cap PD_m.
\]
We use the notation \(PD_m\) to denote the set of all \(m \times m\) symmetric positive definite matrices.

Example (Gaussian contraction axiom). 
Let \(\mathcal{C} = \{ 1 \perp\!\!\!\perp 2 \mid 3,\, 2 \perp\!\!\!\perp 3 \}\). 
By the contraction axiom these two statements imply that \(2 \perp\!\!\!\perp \{ 1, 3 \}\). 
Let us consider this in the Gaussian case. We have
\[
\begin{aligned}
J_{\mathcal{C}} &= \langle \det \Sigma_{\{1,3\}, \{2,3\}},\, \sigma_{23} \rangle \\
    &= \langle \sigma_{12} \sigma_{33} - \sigma_{13} \sigma_{23},\, \sigma_{23} \rangle \\
    &= \langle \sigma_{12} \sigma_{33},\, \sigma_{23} \rangle \\
    &= \langle \sigma_{12},\, \sigma_{23} \rangle \cap \langle \sigma_{33},\, \sigma_{23} \rangle.
\end{aligned}
\]
(which is a primary decomposition).

In particular, the variety defined by the ideal \(J_{\mathcal{C}}\) has two components. 
However, only one of these components is interesting from the probabilistic standpoint, since
\(V(\langle \sigma_{33}, \sigma_{23} \rangle)\) does not intersect the cone of positive definite matrices.

Similar approaches for discrete multivariate distributions via their independence ideal.
This allows us to parameterize families of distributions sometimes more easily, which satisfy certain independence statements.
